% ============================================================================
% Multiplication Table — Grade 7 Math Workbook
% Same data as the practice-tests version but restyled with the workbook's
% angular dashed-callout design language.
% ============================================================================

\thispagestyle{empty}

{
\sffamily

% --- Title ---
\wbSectionHead{\faTh\quad Multiplication Table}{A handy reference — use it whenever you need it!}

% --- Multiplication Table ---
\begin{center}
\renewcommand{\arraystretch}{1.6}
\setlength{\tabcolsep}{0pt}
\fontsize{11}{13}\selectfont

\begin{tabular}{|>{\centering\arraybackslash}m{1.26cm}|>{\centering\arraybackslash}m{1.26cm}|>{\centering\arraybackslash}m{1.26cm}|>{\centering\arraybackslash}m{1.26cm}|>{\centering\arraybackslash}m{1.26cm}|>{\centering\arraybackslash}m{1.26cm}|>{\centering\arraybackslash}m{1.26cm}|>{\centering\arraybackslash}m{1.26cm}|>{\centering\arraybackslash}m{1.26cm}|>{\centering\arraybackslash}m{1.26cm}|>{\centering\arraybackslash}m{1.26cm}|>{\centering\arraybackslash}m{1.26cm}|}
\hline
\rowcolor{funBlue!85}
\textcolor{white}{\bfseries $\times$} &
\textcolor{white}{\bfseries 1} &
\textcolor{white}{\bfseries 2} &
\textcolor{white}{\bfseries 3} &
\textcolor{white}{\bfseries 4} &
\textcolor{white}{\bfseries 5} &
\textcolor{white}{\bfseries 6} &
\textcolor{white}{\bfseries 7} &
\textcolor{white}{\bfseries 8} &
\textcolor{white}{\bfseries 9} &
\textcolor{white}{\bfseries 10} &
\textcolor{white}{\bfseries 11} \\
\hline

\rowcolor{funBlueLight}
\cellcolor{funBlue!85}\textcolor{white}{\bfseries 1} &
1 & 2 & 3 & 4 & 5 & 6 & 7 & 8 & 9 & 10 & 11 \\
\hline

\rowcolor{white}
\cellcolor{funBlue!85}\textcolor{white}{\bfseries 2} &
2 & 4 & 6 & 8 & 10 & 12 & 14 & 16 & 18 & 20 & 22 \\
\hline

\rowcolor{funBlueLight}
\cellcolor{funBlue!85}\textcolor{white}{\bfseries 3} &
3 & 6 & 9 & 12 & 15 & 18 & 21 & 24 & 27 & 30 & 33 \\
\hline

\rowcolor{white}
\cellcolor{funBlue!85}\textcolor{white}{\bfseries 4} &
4 & 8 & 12 & 16 & 20 & 24 & 28 & 32 & 36 & 40 & 44 \\
\hline

\rowcolor{funBlueLight}
\cellcolor{funBlue!85}\textcolor{white}{\bfseries 5} &
5 & 10 & 15 & 20 & 25 & 30 & 35 & 40 & 45 & 50 & 55 \\
\hline

\rowcolor{white}
\cellcolor{funBlue!85}\textcolor{white}{\bfseries 6} &
6 & 12 & 18 & 24 & 30 & 36 & 42 & 48 & 54 & 60 & 66 \\
\hline

\rowcolor{funBlueLight}
\cellcolor{funBlue!85}\textcolor{white}{\bfseries 7} &
7 & 14 & 21 & 28 & 35 & 42 & 49 & 56 & 63 & 70 & 77 \\
\hline

\rowcolor{white}
\cellcolor{funBlue!85}\textcolor{white}{\bfseries 8} &
8 & 16 & 24 & 32 & 40 & 48 & 56 & 64 & 72 & 80 & 88 \\
\hline

\rowcolor{funBlueLight}
\cellcolor{funBlue!85}\textcolor{white}{\bfseries 9} &
9 & 18 & 27 & 36 & 45 & 54 & 63 & 72 & 81 & 90 & 99 \\
\hline

\rowcolor{white}
\cellcolor{funBlue!85}\textcolor{white}{\bfseries 10} &
10 & 20 & 30 & 40 & 50 & 60 & 70 & 80 & 90 & 100 & 110 \\
\hline

\rowcolor{funBlueLight}
\cellcolor{funBlue!85}\textcolor{white}{\bfseries 11} &
11 & 22 & 33 & 44 & 55 & 66 & 77 & 88 & 99 & 110 & 121 \\
\hline

\end{tabular}
\end{center}

\vspace{5mm}

% --- How to Use ---
\begin{wbCallout}{funGreen}
    {\color{funGreen}\faLightbulb}\;\;{\bfseries\color{black!80}How to Use This Table}
    \par\vspace{3pt}
    \small
    To find $\mathbf{4 \times 7}$:
    \begin{enumerate}[leftmargin=*, itemsep=1pt]
        \item Find \textbf{4} in the left column (blue).
        \item Find \textbf{7} in the top row (blue).
        \item Follow the row and column until they meet: the answer is $\mathbf{28}$!
    \end{enumerate}
\end{wbCallout}

\vspace{3mm}

% --- Division Tip ---
\begin{wbCallout}{funOrange}
    {\color{funOrange}\faInfoCircle}\;\;{\bfseries\color{black!80}Division Tip}
    \par\vspace{3pt}
    \small
    You can also use this table for \textbf{division}!
    If you know $28 \div 4 = ?$, find $28$ in the $4$'s row.
    The column header gives you the answer: $\mathbf{7}$!
\end{wbCallout}

}

\newpage
